% -----------------------------
% 執筆にあたっての留意点
% -----------------------------
%
% 研究業績の節には，本研究に関連して発表した業績を列挙します．
% 卒業論文・修士論文で求められることが多い一方，学会に投稿する論文では通常不要です
% （投稿先によっては，匿名審査のために自己引用の書き方に注意が必要です）．
%
% ■ 書き方
% * 種別ごとに分けて書きます（論文誌，査読付き国際会議，査読付き国内会議，
%   研究会・全国大会などの口頭発表，ポスター・デモ発表，受賞 など）
% * 各業績には査読の有無を明記します
% * 著者名，タイトル，発表先の正式名称，巻号・ページ，発表年を漏れなく書きます
% * 発表先は略称ではなく正式名称で書き，必要に応じて略称を併記します
%   （例：「iiWAS 2020」ではなく「Proceedings of the 22nd International Conference on
%     Information Integration and Web-Based Applications \& Services (iiWAS 2020)」）
%   英語の会議名・論文誌名は，日本語に訳さずそのまま書きます
% * 自分が筆頭著者でない業績も，本研究に関係するものは含めます
% * 投稿中・採録決定済みで未発表のものは，その状態を明記します
%   （例：「採録決定，2026年3月発表予定」「投稿中」）
% * 本論文のどの章がどの業績に対応するのかを一言添えると，読み手（審査員）に親切です
%
% ■ 書式は最後に必ず点検してください
% 著者名の表記（姓名の順，イニシャルの有無），ページ番号のハイフンの種類，年の位置などが
% 業績ごとにばらついていないかを確認します．Webから拾ったBibTeXやコピーした文字列をそのまま貼り付けると，
% 表記が揺れがちです（→ how-to-xx/how-to-write-a-paper/1-学術論文の書き方.md 3.7節，
% how-to-xx/how-to-write-a-paper/2-文書の書き方.md 5.3節「表記のゆらぎ」）．
%
% 以下は例です．

\section*{研究業績}
\begin{enumerate}
\item Masaki Suzuki and Yusuke Yamamoto. Analysis of relationship between confirmation bias and web search behavior. In Proceedings of the 22nd International Conference on Information Integration and Web-Based Applications \& Services (iiWAS 2020), pp. 184–191, 2020 (査読あり).
\item Masaki Suzuki and Yusuke Yamamoto. Characterizing the influence of confirmation bias on web search behavior. Frontiers in Psychology, Vol. 12, pp. 1–11, 2021 (査読あり).
\end{enumerate}
